% =============================================================================
%  CHAPITRE 2 — MÉTHODOLOGIE DE L'ÉTUDE
% =============================================================================
\entete{Chapitre 2 — Méthodologie de l'étude}
\chapter{Méthodologie de l'étude}
\label{chap:methodo}

\introchap

La valeur d'un résultat tient à la rigueur du chemin qui y mène. Ce chapitre
expose donc, avec le détail nécessaire à sa reproduction, la démarche suivie
pour conduire l'étude. Il précise successivement la nature de la recherche
entreprise et l'approche méthodologique retenue (section~\ref{sec:nature}), les
variables et les indicateurs au moyen desquels le phénomène a été mesuré
(section~\ref{sec:variables}), la population concernée et l'échantillon
effectivement étudié (section~\ref{sec:echantillon}), les outils de collecte et
les techniques d'analyse des données (section~\ref{sec:outils}), puis la méthode
d'ingénierie logicielle et la démarche d'industrialisation adoptées pour
construire l'artefact (sections~\ref{sec:2tup} et~\ref{sec:devops}). Il se clôt
sur l'énoncé des difficultés rencontrées et des limites méthodologiques
assumées.

% =============================================================================
\section{Nature de l'étude et approche méthodologique}
\label{sec:nature}

\subsection{Niveau et type de recherche}

Le choix du niveau de recherche détermine la profondeur du traitement
scientifique et, par voie de conséquence, les méthodes admissibles. Quatre
niveaux sont habituellement distingués : perceptuel, appréhensif, compréhensif
et intégrateur.

La présente étude relève du \textbf{niveau intégrateur}. Nous ne nous sommes en
effet pas limité à décrire une situation ni à en expliquer les causes : nous
sommes intervenu sur la réalité observée en concevant, en construisant et en
mettant en service un dispositif destiné à la transformer. Le type d'étude
correspondant est la \textbf{recherche interactive}, également désignée dans la
littérature en systèmes d'information sous le nom de \emph{recherche par la
conception} (Hevner et al., 2004). Elle se caractérise par :

\begin{itemize}
  \item une forte interaction avec les acteurs du terrain — le personnel de la
        Fondation a été associé à chaque étape ;
  \item une démarche participative, les besoins ayant été formulés, corrigés et
        arbitrés avec l'encadreur professionnel ;
  \item la résolution d'un problème réel et non d'un cas d'école ;
  \item le recours au prototypage : chaque fonctionnalité a été construite,
        soumise, éprouvée, puis corrigée.
\end{itemize}

Ce niveau excède ce qui est exigé d'un travail de licence, lequel peut se
limiter au niveau perceptuel ou appréhensif. Ce choix se justifie par la nature
même de la commande : la Fondation n'attendait pas un diagnostic, elle attendait
un outil.

\subsection{Approche méthodologique}

Nous avons retenu une \textbf{approche mixte}, associant les apports qualitatifs
et quantitatifs, pour trois raisons.

\begin{enumerate}
  \item La compréhension du besoin — ce qu'une institution patrimoniale attend
        d'une visite à distance — appelle des données qualitatives : entretiens,
        observation, analyse des documents internes.
  \item L'évaluation du dispositif construit appelle au contraire des données
        quantitatives : temps de réponse, taux de réussite, scores de
        correspondance, coûts mesurés.
  \item Le croisement des deux permet de vérifier qu'un dispositif
        techniquement performant est également jugé utile et utilisable par
        ceux à qui il est destiné, conformément au modèle d'acceptation évoqué
        au chapitre précédent.
\end{enumerate}

Le tableau~\ref{tab:approche} récapitule la répartition des deux approches selon
les objectifs poursuivis.

\begin{table}[htbp]
\centering
\caption{Articulation des approches qualitative et quantitative selon les objectifs}
\label{tab:approche}
\small
\begin{tabularx}{\textwidth}{|L{2.4cm}|L{3.4cm}|X|}
\hline
\rowcolor{keycetete}
\thead{Approche} & \thead{Objectifs concernés} & \thead{Nature des données recueillies} \\ \hline
Qualitative & Compréhension du besoin, étude de l'existant, arbitrages de conception &
Comptes rendus d'entretiens, notes d'observation, documents internes de la Fondation, retours d'usage \\ \hline
Quantitative & Évaluation du dispositif, vérification des hypothèses &
Mesures de performance, scores de correspondance documentaire, taux de réussite, relevés de coûts, réponses au questionnaire \\ \hline
\end{tabularx}
\source{Élaboré par nos soins.}
\end{table}

% =============================================================================
\section{Variables et indicateurs de l'étude}
\label{sec:variables}

\subsection{Identification des variables}

L'hypothèse générale énoncée en introduction affirme qu'une plateforme réunissant
quatre composantes permet de restituer à distance l'essentiel de l'expérience de
visite. La \textbf{variable expliquée} est donc :

\begin{itemize}
  \item \textbf{VE — la qualité de l'expérience de visite à distance}, entendue
        comme la capacité du dispositif à procurer au visiteur éloigné une
        perception des lieux, une compréhension des œuvres et une possibilité
        d'interrogation comparables à celles d'une visite sur place.
\end{itemize}

Quatre \textbf{variables explicatives} composent cette hypothèse, une par
hypothèse spécifique :

\begin{itemize}
  \item \textbf{V1 — la restitution spatiale des espaces et des œuvres} ;
  \item \textbf{V2 — la mise en relation documentaire entre institutions} ;
  \item \textbf{V3 — l'ancrage documentaire de l'assistant conversationnel} ;
  \item \textbf{V4 — le cloisonnement et l'ouverture de l'architecture}.
\end{itemize}

\subsection{Opérationnalisation : les indicateurs}

Une variable n'est exploitable que si l'on sait la mesurer. Le
tableau~\ref{tab:variables} présente, pour chaque variable, les indicateurs
retenus, leur unité et la source de la mesure.

\begin{table}[htbp]
\centering
\caption{Variables, indicateurs et sources de mesure}
\label{tab:variables}
\footnotesize
\begin{tabularx}{\textwidth}{|L{2.5cm}|X|L{2.2cm}|L{2.6cm}|}
\hline
\rowcolor{keycetete}
\thead{Variable} & \thead{Indicateurs} & \thead{Unité} & \thead{Source de la mesure} \\ \hline
\multirow{5}{*}{\parbox{2.5cm}{\textbf{V1}\\ Restitution spatiale}}
 & Nombre de salles restituées en panorama & effectif & Base de données \\ \cline{2-4}
& Nombre de points d'intérêt cliquables & effectif & Base de données \\ \cline{2-4}
& Délai d'affichage d'une salle & secondes & Mesure navigateur \\ \cline{2-4}
& Écart dimensionnel du modèle 3D & centimètres & Mesure sur maillage \\ \cline{2-4}
& Couverture des appareils testés & \% & Essais sur appareils \\ \hline
\multirow{5}{*}{\parbox{2.5cm}{\textbf{V2}\\ Mise en relation documentaire}}
 & Nombre de collections interrogées & effectif & Journal des appels \\ \cline{2-4}
& Nombre d'œuvres apparentées identifiées & effectif & Table des rapprochements \\ \cline{2-4}
& Score de correspondance & /100 & Algorithme de notation \\ \cline{2-4}
& Taux de validation par le conservateur & \% & Table des rapprochements \\ \cline{2-4}
& Nombre de pays de dispersion & effectif & Données de dispersion \\ \hline
\multirow{5}{*}{\parbox{2.5cm}{\textbf{V3}\\ Ancrage documentaire de l'assistant}}
 & Taux de réponses appuyées sur un document interne & \% & Grille d'évaluation \\ \cline{2-4}
& Taux d'aveu d'ignorance en l'absence de source & \% & Grille d'évaluation \\ \cline{2-4}
& Taux de refus des questions hors périmètre & \% & Grille d'évaluation \\ \cline{2-4}
& Délai moyen de réponse & secondes & Journal serveur \\ \cline{2-4}
& Questions restées sans réponse & effectif & Journal des questions \\ \hline
\multirow{4}{*}{\parbox{2.5cm}{\textbf{V4}\\ Cloisonnement et ouverture}}
 & Tables protégées par une règle de sécurité & effectif & Schéma de la base \\ \cline{2-4}
& Tentatives d'accès inter-organisations refusées & \% & Essais d'intrusion \\ \cline{2-4}
& Délai de création d'une organisation & minutes & Essai de bout en bout \\ \cline{2-4}
& Lignes de code à modifier pour un nouvel arrivant & effectif & Analyse du code \\ \hline
\end{tabularx}
\source{Élaboré par nos soins.}
\end{table}

\encadre{Note de méthode}{%
Le choix d'indicateurs \emph{observables sur le système lui-même} — et non
seulement déclarés par les utilisateurs — répond à une exigence de vérifiabilité.
Un taux de refus des questions hors périmètre se constate ; une satisfaction
déclarée s'interprète. Les deux ont été recueillis, mais les premiers fondent la
vérification des hypothèses et les seconds l'éclairent.}

% =============================================================================
\section{Population et échantillon}
\label{sec:echantillon}

\subsection{Population de l'étude}

La population de l'étude est constituée de trois ensembles d'acteurs concernés
par la valorisation numérique du patrimoine de la Fondation Jean Félicien Gacha
au cours de la période du 20~juin au 20~août 2026 :

\begin{itemize}
  \item \textbf{le personnel de la Fondation et du \emph{Think Tank} Foudjem},
        c'est-à-dire les personnes intervenant dans la gestion des espaces, des
        collections, de l'accueil et de la communication ;
  \item \textbf{les visiteurs de la Fondation}, présents sur le site durant la
        période, âgés de dix-huit ans et plus ;
  \item \textbf{les institutions patrimoniales de la région de l'Ouest}
        susceptibles de rejoindre la plateforme : chefferies, musées et
        associations culturelles.
\end{itemize}

\subsection{Technique d'échantillonnage et taille de l'échantillon}

Compte tenu de la taille réduite et bien identifiée des populations concernées,
un échantillonnage probabiliste n'aurait apporté aucun gain de représentativité.
Nous avons donc retenu deux techniques \textbf{non probabilistes} :

\begin{itemize}
  \item \textbf{le choix raisonné} pour le personnel : ont été retenues les
        personnes exerçant effectivement une responsabilité sur les espaces, les
        collections ou la communication, seules à même de renseigner les
        pratiques réelles ;
  \item \textbf{la convenance} pour les visiteurs : ont été sollicitées les
        personnes présentes sur le site durant les journées d'observation et
        ayant accepté de répondre.
\end{itemize}

Le tableau~\ref{tab:echantillon} présente la composition de l'échantillon.

\begin{table}[htbp]
\centering
\caption{Composition de l'échantillon de l'étude}
\label{tab:echantillon}
\small
\begin{tabularx}{\textwidth}{|X|C{2.4cm}|C{2.6cm}|L{3.4cm}|}
\hline
\rowcolor{keycetete}
\thead{Catégorie d'acteurs} & \thead{Effectif} & \thead{Technique} & \thead{Outil de collecte} \\ \hline
Personnel de la Fondation et du \emph{Think Tank} Foudjem & 06 & Choix raisonné & Entretien semi-directif \\ \hline
Visiteurs présents sur le site & 42 & Convenance & Questionnaire \\ \hline
Institutions patrimoniales approchées & 03 & Choix raisonné & Entretien exploratoire \\ \hline
\rowcolor{keycetotal}
\textbf{Total} & \textbf{51} & — & — \\ \hline
\end{tabularx}
\source{Enquête de terrain, juin--août 2026.}
\end{table}

Cet échantillon ne prétend pas à la représentativité statistique. Il vise la
\emph{saturation de l'information} : les entretiens ont été conduits jusqu'à ce
que les échanges cessent d'apporter des éléments nouveaux, critère usuel en
recherche qualitative.

% =============================================================================
\section{Outils de collecte et techniques d'analyse des données}
\label{sec:outils}

\subsection{Outils de collecte}

Quatre outils ont été mobilisés, chacun rattaché à un type de donnée précis.

\begin{enumerate}
  \item \textbf{Le guide d'entretien semi-directif} (annexe~\ref{ann:guide}).
        Organisé autour de quatre thèmes — pratiques actuelles de valorisation,
        difficultés rencontrées, attentes vis-à-vis d'un dispositif à distance,
        contraintes de moyens — il a permis de recueillir les besoins auprès du
        personnel. Sa forme semi-directive laisse à l'interlocuteur la liberté
        d'aborder des points non anticipés, ce qui s'est révélé décisif : la
        question de la dispersion des œuvres apparentées n'était pas prévue au
        départ et a émergé d'un entretien.

  \item \textbf{Le questionnaire} (annexe~\ref{ann:questionnaire}). Adressé aux
        visiteurs, il comporte dix-huit questions fermées et deux questions
        ouvertes, réparties en quatre rubriques : profil du répondant, pratiques
        numériques, attentes relatives à une visite à distance, appréciation du
        prototype après démonstration. Un pré-test auprès de cinq personnes a
        conduit à reformuler trois questions jugées ambiguës et à supprimer une
        question redondante.

  \item \textbf{La grille d'observation.} Utilisée lors des visites guidées,
        elle a servi à relever le parcours effectif des visiteurs dans les
        espaces, les points d'arrêt, les questions spontanément posées au guide
        et la durée moyenne de stationnement devant les œuvres. Ces relevés ont
        directement orienté la conception du parcours immersif et
        l'identification des questions que l'assistant devait savoir traiter.

  \item \textbf{La fiche documentaire.} Elle a encadré le dépouillement des
        documents internes de la Fondation — inventaires, notices, plaquettes,
        photographies d'archives — ainsi que la revue de la littérature
        scientifique et de la documentation technique des éditeurs.
\end{enumerate}

\subsection{Techniques d'analyse}

\subsubsection{Analyse qualitative}

Les entretiens ont fait l'objet d'une \textbf{analyse thématique de contenu}.
Après retranscription, les propos ont été découpés en unités de sens, codées,
puis regroupées en catégories. Cinq thèmes ont émergé : l'accessibilité, la
perception de l'espace, la mise en relation des objets, la fiabilité de
l'information et la pérennité économique du dispositif. Aucun logiciel spécialisé
n'a été employé : le volume traité — six entretiens — ne le justifiait pas, et un
dépouillement manuel systématique offrait une maîtrise supérieure du codage.

\subsubsection{Analyse quantitative}

Les réponses au questionnaire ont été traitées par \textbf{statistique
descriptive} : effectifs, fréquences, pourcentages et moyennes, présentés sous
forme de tableaux et de graphiques. Les mesures relevées sur le système
— délais, scores, taux — ont fait l'objet de comparaisons avant/après, notamment
pour évaluer l'apport de l'enrichissement des requêtes documentaires. Le tableur
a suffi à ces traitements.

\subsubsection{Évaluation de l'assistant conversationnel}

L'assistant a été évalué au moyen d'une \textbf{grille d'évaluation} appliquée à
un jeu de questions construit à l'avance
(annexe~\ref{ann:grille}). Ce jeu comporte quatre familles de questions :

\begin{itemize}
  \item des questions \textbf{documentées}, dont la réponse figure dans les
        textes fournis par la Fondation — l'assistant doit répondre exactement ;
  \item des questions \textbf{non documentées}, portant sur le périmètre de la
        Fondation mais sans source disponible — l'assistant doit reconnaître
        qu'il ne sait pas ;
  \item des questions \textbf{hors périmètre} (météo, actualité, informatique
        générale) — l'assistant doit refuser poliment ;
  \item des questions \textbf{pièges}, formulées de manière à suggérer une
        réponse fausse — l'assistant ne doit pas suivre la suggestion.
\end{itemize}

Cette construction est délibérée : elle mesure non seulement ce que l'assistant
sait, mais surtout ce qu'il refuse d'inventer, qui est le critère décisif en
médiation culturelle.

% =============================================================================
\section{Méthode d'ingénierie logicielle}
\label{sec:2tup}

\subsection{Choix du langage de modélisation}

Le langage de modélisation unifié (UML) a été retenu pour la conception. Ce
choix se justifie par trois arguments :

\begin{itemize}
  \item \textbf{la standardisation} : UML étant une norme reconnue de l'Object
        Management Group, les modèles produits restent lisibles par tout
        intervenant ultérieur, ce qui importe pour un système destiné à être
        repris par la Fondation ;
  \item \textbf{la souplesse} : ses diagrammes couvrent aussi bien la structure
        des données que les flux d'interaction, ce qui convient à un système
        composite associant gestion, immersion et assistance ;
  \item \textbf{la qualité de la documentation produite} : les diagrammes
        constituent une documentation durable, indispensable à la maintenance.
\end{itemize}

UML~2 propose treize diagrammes, répartis entre une vue statique et une vue
dynamique, comme le montre le tableau~\ref{tab:uml}. Parmi eux, trois ont été
retenus : le diagramme de cas d'utilisation, le diagramme de classes et le
diagramme de séquence, complétés par un diagramme de déploiement pour
l'infrastructure.

\begin{table}[htbp]
\centering
\caption{Les deux vues d'UML~2 et leurs diagrammes respectifs}
\label{tab:uml}
\small
\begin{tabularx}{\textwidth}{|X|X|}
\hline
\rowcolor{keycetete}
\thead{Diagrammes structurels (vue statique)} & \thead{Diagrammes comportementaux (vue dynamique)} \\ \hline
Diagramme de classes & Diagramme de cas d'utilisation \\ \hline
Diagramme d'objets & Diagramme de séquence \\ \hline
Diagramme de composants & Diagramme d'activités \\ \hline
Diagramme de déploiement & Diagramme d'états-transitions \\ \hline
Diagramme de paquetages & Diagramme de communication \\ \hline
Diagramme de structure composite & Diagramme global d'interaction \\ \hline
& Diagramme de temps \\ \hline
\end{tabularx}
\source{Adapté de la spécification UML~2 de l'Object Management Group.}
\end{table}

\subsection{Choix du processus : 2TUP}

UML est un langage, non une méthode : il faut lui adjoindre un processus. Parmi
les processus fondés sur UML — RUP, OMT, OOSE, 2TUP — notre choix s'est porté sur
le \textbf{2TUP} (\emph{Two Track Unified Process}), déclinaison du processus
unifié développée pour répondre aux systèmes d'information soumis à des
changements constants.

L'expression « deux branches » désigne les deux axes que le processus traite
séparément avant de les réunir :

\begin{itemize}
  \item la \textbf{branche fonctionnelle}, qui part de l'expression des besoins
        des utilisateurs et aboutit au modèle du métier, indépendamment de toute
        technique ;
  \item la \textbf{branche technique}, qui part des contraintes techniques
        — matériel, réseau, plateformes cibles, sécurité — et aboutit à
        l'architecture ;
  \item la \textbf{branche de réalisation}, qui fusionne les deux précédentes en
        une conception préliminaire puis détaillée, avant le codage et les
        essais.
\end{itemize}

Ce processus s'est révélé particulièrement adapté à notre cas. Les contraintes
techniques du projet — impossibilité d'ajouter des dépendances logicielles,
nécessité d'un fonctionnement sur téléphone d'entrée de gamme, exigence de
cloisonnement entre organisations — étaient si structurantes qu'elles méritaient
d'être instruites en parallèle des besoins, et non après eux.

\begin{figure}[htbp]
\centering
\aplacer{7cm}{Représenter le cycle en Y du processus 2TUP : à gauche la branche
fonctionnelle (capture des besoins fonctionnels $\rightarrow$ analyse), à droite
la branche technique (capture des besoins techniques $\rightarrow$ conception
générique), la jonction en conception préliminaire, puis conception détaillée,
codage et essais, enfin recette. Faire figurer, sur chaque branche, les
livrables produits dans le cadre du présent projet.}
\caption{Le cycle en Y du processus 2TUP appliqué au projet}
\label{fig:2tup}
\source{Adapté de Roques et Vallée (2007).}
\end{figure}

% =============================================================================
\section{Démarche d'industrialisation : le choix du DevOps}
\label{sec:devops}

\subsection{Justification du choix}

La construction d'un logiciel ne s'achève pas à l'écriture du code : encore
faut-il le livrer, le corriger et le faire évoluer sans interrompre le service.
Trois raisons ont conduit à adopter une démarche DevOps dès le début du projet
plutôt qu'à la fin.

\begin{enumerate}
  \item \textbf{La sécurité du travail.} Un stage de deux mois ne laisse aucune
        marge pour reconstruire ce qui aurait été cassé sans qu'on s'en
        aperçoive. Des contrôles automatiques déclenchés à chaque modification
        signalent la régression en quelques minutes.

  \item \textbf{La reproductibilité.} La Fondation devra un jour reprendre le
        système. Une mise en production qui repose sur une suite de gestes
        manuels mémorisés par une seule personne n'est pas transmissible ; une
        mise en production décrite dans un fichier l'est.

  \item \textbf{La sensibilité du dispositif aux détails d'exploitation.}
        L'expérience du projet a montré qu'un modèle tridimensionnel servi avec
        une étiquette de type erronée s'affiche correctement dans le navigateur
        mais empêche toute session de réalité augmentée. Ce type d'anomalie ne
        se voit pas à l'œil : il se vérifie automatiquement, ou il passe
        inaperçu jusqu'au jour de la démonstration.
\end{enumerate}

\subsection{Les pratiques retenues}

Quatre pratiques ont été mises en œuvre :

\begin{itemize}
  \item \textbf{le contrôle de version} de l'intégralité du projet, code
        applicatif et description de l'infrastructure comprise ;
  \item \textbf{l'intégration continue} : chaque dépôt de modification déclenche
        la compilation et une batterie de contrôles portant sur le résultat
        produit ;
  \item \textbf{le déploiement continu} : une modification acceptée est publiée
        automatiquement, avec purge du cache de diffusion et vérification, après
        coup, de ce que reçoit réellement le visiteur ;
  \item \textbf{l'infrastructure décrite par le code}, afin que l'environnement
        de production puisse être recréé à l'identique.
\end{itemize}

La chaîne effectivement mise en place et son évolution proposée sont décrites
respectivement au chapitre~3 et au chapitre~4.

% =============================================================================
\section{Outils et environnement de l'étude}
\label{sec:outilsetude}

Le tableau~\ref{tab:outils} recense les outils employés, en distinguant leur
rôle dans le travail.

\begin{table}[htbp]
\centering
\caption{Outils, technologies et environnements mobilisés}
\label{tab:outils}
\footnotesize
\begin{tabularx}{\textwidth}{|L{3.0cm}|L{3.9cm}|X|}
\hline
\rowcolor{keycetete}
\thead{Domaine} & \thead{Outil ou technologie} & \thead{Rôle dans le travail} \\ \hline
\multirow{4}{*}{\parbox{3.0cm}{Interface}}
 & Vue~3 & Cadriciel de construction de l'interface \\ \cline{2-3}
& Vite~5 & Chaîne de compilation et serveur de développement \\ \cline{2-3}
& PrimeVue~4 & Bibliothèque de composants d'interface pour l'ERP \\ \cline{2-3}
& Pinia, Vue~Router, Vue~I18n & État applicatif, navigation, bilinguisme \\ \hline
\multirow{4}{*}{\parbox{3.0cm}{Données et services}}
 & PostgreSQL~17 & Base de données relationnelle \\ \cline{2-3}
& Supabase & Authentification, sécurité au niveau de la ligne, fonctions serveur \\ \cline{2-3}
& \texttt{pgvector} (index HNSW) & Stockage et recherche des plongements \\ \cline{2-3}
& Deno & Exécution des fonctions serveur \\ \hline
\multirow{3}{*}{\parbox{3.0cm}{Immersion et 3D}}
 & WebGL (moteur écrit sur mesure) & Rendu des panoramas à 360~degrés et du relief \\ \cline{2-3}
& \texttt{model-viewer}, WebXR & Affichage 3D et réalité augmentée \\ \cline{2-3}
& Meshroom, Blender & Photogrammétrie et préparation des maillages \\ \hline
\multirow{3}{*}{\parbox{3.0cm}{Intelligence artificielle}}
 & Groq (modèles Llama) & Génération et reformulation de texte \\ \cline{2-3}
& \texttt{gte-small} & Production des plongements lexicaux \\ \cline{2-3}
& ElevenLabs & Synthèse vocale du guide \\ \hline
\multirow{5}{*}{\parbox{3.0cm}{Industrialisation}}
 & Git, GitHub & Contrôle de version et hébergement du dépôt \\ \cline{2-3}
& GitHub Actions & Intégration et déploiement continus \\ \cline{2-3}
& Docker & Conteneurisation de l'application \\ \cline{2-3}
& Terraform & Description de l'infrastructure \\ \cline{2-3}
& AWS (S3, CloudFront, Route~53, ACM, IAM) & Hébergement et diffusion \\ \hline
\multirow{3}{*}{\parbox{3.0cm}{Conception et rédaction}}
 & Visual Studio Code & Édition du code \\ \cline{2-3}
& \texttt{draw.io} & Réalisation des schémas \\ \cline{2-3}
& \LaTeX & Composition du présent rapport \\ \hline
\end{tabularx}
\source{Élaboré par nos soins.}
\end{table}

% =============================================================================
\section{Difficultés rencontrées et limites méthodologiques}

L'honnêteté méthodologique commande d'exposer les obstacles rencontrés, car ils
conditionnent la portée des résultats.

\begin{enumerate}
  \item \textbf{Impossibilité d'installer de nouvelles bibliothèques.} Le poste
        de travail utilisé pendant le stage ne permettait pas l'installation de
        nouvelles dépendances logicielles. Cette contrainte, subie et non
        choisie, a eu une conséquence majeure : le moteur de panorama, le moteur
        de relief, l'encodeur de codes QR, le générateur de fichiers 3D et le
        globe de dispersion ont dû être écrits intégralement. Ce qui semblait
        une entrave s'est révélé instructif — le dispositif final ne dépend
        d'aucune bibliothèque tierce pour ses fonctions les plus visibles — mais
        a coûté un temps considérable.

  \item \textbf{Absence de clés d'accès à quatre grandes collections.} Europeana,
        le Rijksmuseum, la Smithsonian Institution et les Harvard Art Museums
        exigent une clé que nous n'avons pu obtenir dans les délais. Le
        rapprochement documentaire porte donc sur cinq sources et non sur neuf,
        ce qui minore mécaniquement le nombre d'œuvres apparentées trouvées.

  \item \textbf{Comparaison portant sur le texte et non sur l'image.} Faute de
        processeur graphique et de possibilité d'installer un modèle de vision,
        les plongements ont été calculés sur les descriptions textuelles des
        œuvres et non sur leurs photographies. Deux objets visuellement
        semblables mais décrits différemment peuvent donc échapper au
        rapprochement. Cette limite est explicitement assumée et figure parmi
        les perspectives.

  \item \textbf{Essais de réalité augmentée limités.} Le décodage des codes QR
        par caméra et le démarrage effectif d'une session de réalité augmentée
        n'ont pu être éprouvés que sur un nombre restreint d'appareils, le poste
        de développement ne disposant pas des interfaces nécessaires.

  \item \textbf{Taille de l'échantillon.} Conduite auprès d'une institution
        unique et d'un échantillon de convenance, l'étude ne permet pas de
        généralisation statistique. Ses enseignements sont transférables sous
        réserve de vérification.
\end{enumerate}

\conclchap

Ce chapitre a exposé la démarche de l'étude : une recherche interactive de
niveau intégrateur, conduite selon une approche mixte, dont les quatre variables
explicatives ont été rendues mesurables par un ensemble d'indicateurs
observables sur le système lui-même. Il a précisé la population, l'échantillon,
les outils de collecte et les techniques d'analyse, puis la méthode
d'ingénierie — le processus 2TUP appuyé sur UML — et la démarche
d'industrialisation retenue. Les limites assumées y ont été énoncées sans
détour, car elles bornent la portée des résultats. Le chapitre suivant présente
le terrain, les données recueillies et les résultats obtenus.
